\section{问题、结论与证明路线}\label{sec:introduction}

相邻整数幂之间的素数问题问：固定整数 $k\ge2$，是否对每个正整数 $n$，开区间 $(n^k,(n+1)^k)$ 中都存在素数？$k=2$ 是通常的勒让德猜想。本文处理 $k=6$，并将每一个外部数学输入固定为可核对的精确陈述。

这一问题有两个尺度。对大 $n$，区间长度约为 $kn^{k-1}$，可以借助零点分布与显式公式控制素数计数；对小 $n$，渐近估计常数较大，直接素性见证和显式短区间结果更有效。两部分之间的有限范围必须被连续覆盖，而不能由若干采样点的成功代替。

\begin{theorem}[主定理]\label{thm:main}
接受第\ref{sec:inputs}节的 H1--H7。对每个整数 $n\ge1$，存在素数 $p$，使
\begin{equation}\label{eq:main}
n^6<p<n^6+6n^5<(n+1)^6.
\end{equation}
\end{theorem}

主定理先证明较短区间内的素数存在性。右侧第二个严格不等式完全来自二项式展开；数论工作集中在第一个区间。对于正实数 $x$，记
\[
\mathcal S_6(x):\quad \exists p\text{ prime},\quad x<p<x+6x^{5/6},
\qquad
\mathcal S_4(x):\quad \exists p\text{ prime},\quad x<p<x+4x^{3/4}.
\]
当 $x\ge1$ 时，$4x^{3/4}<6x^{5/6}$，所以 $\mathcal S_4(x)$ 足以推出 $\mathcal S_6(x)$。这一长度比较允许不同范围采用各自有效的参数。

\subsection{研究前提与目标的选择}
本项目选择接受 QRH$(7/8)$，即实际黎曼 $\zeta$ 函数每个非平凡零点的实部至多为 $7/8$。研究动机来自 OpenAI 公开的相关论文与专项形式化材料\cite{OpenAIQRH}；本文取其明确的零点位置结论作为 H1。显式密度输入另列为 H2，其他五项文献结果列为 H3--H7。这样一来，主定理的数学含义由这些陈述本身固定，上游证据的更新可以独立记录。

路线的阶段研究目标仍为 $k=4$，长期目标为 $k=3$。$k=6$ 是当前已经闭合全部范围、并完成下游形式化的成果。本文会解释哪些估计可复用到较小指数，但只把已经闭合的范围写成定理。

\subsection{覆盖结构}
令 $L=\log x$。最终证明按下表组织；两个相邻范围的共同边界都有明确归属。
\begin{center}\small
\begin{tabular}{p{0.26\linewidth}p{0.25\linewidth}p{0.36\linewidth}}
\toprule
范围 & 局部目标 & 主要方法\\
\midrule
整数 $1\le n\le1242317999994$ & $n^6<p<n^6+6n^5$ & 有限素性见证、CHL 十二行和整数阈值\\
$167\le L\le715/4$ & $\mathcal S_6(x)$ & 有限高度 RH、缩放核及反向密度比较\\
$715/4\le L\le400$ & $\mathcal S_6(x)$ & 困难基准证书及连续传播\\
$400\le L\le700$ & $\mathcal S_4(x)$ & 远段梯形证书及连续传播\\
$700\le L$ & $\mathcal S_4(x)$ & 共同截断显式公式与解析尾界\\
\bottomrule
\end{tabular}\end{center}
整数下段与 $L=167$ 处有已证明的重叠。最终拼接只需将 $x=n^6$ 代入相应实变量结论，再作严格的长度比较。

\subsection{本项目的具体工作}
本文所组织的证明包含四项相互连接的工作。首先，用实际 $\zeta$ 零点及解析重数建立加权显式公式，并把素数幂的贡献从正的加权和中扣除。其次，将零点密度估计变成非负有限网格预算，并处理无限高度中所有残余权重。再次，用指数矩的低秩分解产生可由内核检查的整数证书，再证明证书随实参数的连续传播。最后，把这些估计与有限整数段、共同截断尾部组合成全范围结论。

证明的三个代表性完整预算为 $0.939772$、$0.966283$ 和 $0.450022$，均严格小于归一化主项 $1$。它们分别服务于困难基准区间、下中段和远段。各预算的分项及推导在后文给出；这些数字不是未经证明的程序输出前提。

\subsection{阅读路径}
第\ref{sec:inputs}节固定输入，第\ref{sec:preliminaries}节统一数学对象。第\ref{sec:trapezoid}--\ref{sec:coverage}节建立有限解析段；第\ref{sec:tail}节处理无限尾部；第\ref{sec:lower}节处理整数下段；第\ref{sec:assembly}节完成主定理。后两节讨论方法适用性和可复现形式化。关键命题后的形式化对应列出文件与声明名，方便读者在数学证明和源码之间往返。

\leanref{FinalResult}{sixth_powers_from_named_inputs; sixth_power_short_from_named_inputs}
